% Spelersuitleg bij de drankje-inzet (#1004), verteld door Coach Rudy.
%
% Bouwen:  ./scripts/build-uitleg-pdf.sh
% Dat script zet de webp-avatars uit src/features/coach/components/rudi_avatars
% om naar png en roept daarna xelatex aan — LaTeX leest zelf geen webp.
%
% Bewust zonder emoji: XeLaTeX kan de kleuren-emojifont van macOS niet
% insluiten. De kleuren spiegelen de app: --dorst (amber van de traktatie),
% --coach (Rudy's warme amber) en --accent (smaragd) uit src/app/index.css.
\documentclass[10pt,a4paper]{article}

\usepackage{fontspec}
\usepackage[dutch]{babel}
\usepackage[margin=1.2cm]{geometry}
\usepackage{xcolor}
\usepackage{tcolorbox}
\usepackage{enumitem}
\usepackage{multicol}
\usepackage{parskip}
\setlength{\parskip}{4pt}
\usepackage{tikz}
\usepackage{graphicx}

\tcbuselibrary{skins}

\setmainfont{Avenir Next}
% Iets strakkere regelval: zo past de hele uitleg op één blad.
\linespread{0.97}

% Kleuren uit src/app/index.css — dezelfde families als in de app zelf.
\definecolor{dorst}{HTML}{9A6010}      % --dorst: de traktatie
\definecolor{dorstsoft}{HTML}{FDF2E0}  % --dorst-soft
\definecolor{dorstline}{HTML}{ECD4A8}  % --dorst-line
\definecolor{coachsoft}{HTML}{FDF0DC}  % --coach-soft
\definecolor{coachline}{HTML}{F0D3A3}  % --coach-line
\definecolor{coachdiep}{HTML}{5A3410}  % --coach-diep
\definecolor{inkt}{HTML}{1A2620}       % --ink
\definecolor{inktzacht}{HTML}{586A60}  % --ink-soft-strong
\definecolor{smaragd}{HTML}{0C8A5F}    % --accent

\color{inkt}
\pagestyle{empty}
\setlist[itemize]{leftmargin=1.2em,itemsep=2.5pt,topsep=3pt}

\newcommand{\kop}[1]{%
  \vspace{0.35em}%
  {\color{dorst}\large\bfseries #1}\par\vspace{0.1em}%
}

% Rudy's ronde avatar. De bron is vierkant met zwarte hoeken, dus wegknippen
% tot een cirkel — net als .coach-avatar in de app. #2 is de diameter; de
% straal is de helft daarvan, vandaar \dimexpr en niet ,,2#2'' (dat plakt de
% cijfers aan elkaar tot 21cm).
\newcommand{\rudikop}[2][rudi-portret]{%
  \begin{tikzpicture}
    \clip (0,0) circle (\dimexpr#2/2\relax);
    \node[inner sep=0pt] at (0,0) {\includegraphics[width=#2]{#1}};
  \end{tikzpicture}%
}

% Eén sneer van de coach: avatar links, tekstballon rechts, in het coach-amber.
\newcommand{\rudi}[2][rudi-portret]{%
  \begin{tcolorbox}[colback=coachsoft,colframe=coachline,boxrule=1pt,arc=6pt,
    left=8pt,right=10pt,top=5pt,bottom=5pt]
    \noindent
    \begin{minipage}[c]{1.3cm}\rudikop[#1]{1.15cm}\end{minipage}%
    \begin{minipage}[c]{\dimexpr\linewidth-1.4cm\relax}
      {\color{coachdiep}\itshape #2}
    \end{minipage}
  \end{tcolorbox}%
}

\begin{document}

% --- Titelblok -------------------------------------------------------------
\begin{tcolorbox}[colback=dorstsoft,colframe=dorstline,boxrule=1pt,
  arc=6pt,left=12pt,right=12pt,top=7pt,bottom=7pt]
  \noindent
  \begin{minipage}[c]{\dimexpr\linewidth-2.6cm\relax}
    {\color{dorst}\LARGE\bfseries Nieuw: de drankje-inzet}\par\vspace{3pt}
    {\color{inktzacht}Vanaf nu speelt u niet alleen voor uw rating, maar ook
    voor een pint. De verliezers trakteren.}
  \end{minipage}%
  \hfill
  \begin{minipage}[c]{2.3cm}
    \hfill\rudikop{2.1cm}
  \end{minipage}
\end{tcolorbox}

\rudi{Ik heb veel meegemaakt. Elf man die niet wisten waar ze moesten staan. Een
watersproeier die mij persoonlijk viseerde. Maar een club die maandenlang
\emph{om de eer} speelt --- dat is het grootste tactische vacu\"um dat ooit in
mijn notitieboekje belandde. Vanaf vandaag ligt er iets op tafel. Een glas, om
precies te zijn.}

% --- Wat is het ------------------------------------------------------------
\kop{Waar gaat dit over?}

Vroeger eindigde een match met een cijfer. Nu eindigt hij met een
\textbf{rekening}. Bij het plannen van een wedstrijd kiest u een drankje van de
Belgische kaart --- van een Jupiler tot een Tripel Karmeliet, of gewoon een Spa
plat voor wie de dag erna nog wil kunnen bewegen. Wie verliest, trakteert. De
app zet de afspraak op de wedstrijdkaart en in de feed, en houdt op uw profiel
bij hoeveel u al hebt weggegeven. Betalen doet u zelf.

% --- Hoe werkt het ---------------------------------------------------------
\kop{In drie stappen}

\begin{enumerate}[leftmargin=1.4em,itemsep=3pt,topsep=3pt]
  \item \textbf{Plan uw match} zoals altijd. In stap 2, onder het baantype,
        staat nu \emph{Drankje-inzet}. Zoeken, aantikken, klaar.
  \item \textbf{Kies het aantal.} Standaard \'e\'en consumptie per winnaar.
        Voelt u zich sterk? Zet er twee op. Of drie. Tot tien --- maar dan
        hebben we het niet meer over padel.
  \item \textbf{Speel de pot.} Na afloop staat op de match wie wat schuldig is.
        Aan de bar tikt u op \textbf{,,Traktatie ingelost''} en de rekening is
        vereffend.
\end{enumerate}

\rudi[rudi-gemeen]{Stap twee is waar de meesten sneuvelen. Ik zie ze denken:
,,drie Duvels, dat maakt indruk''. Wat het maakt, is een dure avond voor iemand
die zijn tweede opslag niet vertrouwt. Zet in naar uw niveau, niet naar uw
zelfbeeld.}

% --- De kaart --------------------------------------------------------------
\kop{Wat staat er op de kaart?}

\begin{multicols}{2}
\textbf{25 Belgische bieren} --- van de klassiekers (Jupiler, Stella, Maes) over
de trappisten (Westmalle, Chimay, Orval, Rochefort) tot het zware werk
(Delirium Tremens, Troubadour Magma). Er zit gegarandeerd iets tussen dat uw
tegenstander niet graag afrekent.

\columnbreak

\textbf{10 frisdranken en waters} --- Cola, Fanta, Sprite, Ice Tea, Spa,
Chaudfontaine, Perrier, Red Bull, Schweppes en Oasis. Even geldig als inzet, en
u speelt de week erna beter.
\end{multicols}

% --- De regels -------------------------------------------------------------
\kop{De spelregels}

\begin{tcolorbox}[colback=white,colframe=dorstline,boxrule=1pt,arc=6pt,
  left=11pt,right=11pt,top=6pt,bottom=6pt]
\begin{itemize}
  \item \textbf{Tot de aftrap kunt u nog wijzigen.} Daarna ligt de inzet vast.
        Geen enkele verliezer draait een Westmalle achteraf nog om in een
        kraantjeswater.
  \item \textbf{Gelijkspel? Inzet vervallen.} Niemand wint, dus niemand
        trakteert. Iedereen betaalt zijn eigen glas.
  \item \textbf{Per winnaar.} Verliest u een dubbelspel met \'e\'en drankje
        als inzet, dan staat u er twee te betalen. Reken even door voor u stoer
        doet.
  \item \textbf{Iedereen die meespeelt mag de inzet zetten}, ook op een ronde
        die de app zelf inplande --- en afvinken mag ook iedereen.
  \item \textbf{Uw rating blijft erbuiten.} Dit is een weddenschap tussen
        vrienden, geen tweede lef-tip. Alleen uw eer (en uw portefeuille) staat
        op het spel.
\end{itemize}
\end{tcolorbox}

% --- Profiel ---------------------------------------------------------------
\kop{En op uw profiel?}

Onder \emph{,,Aan de bar''} staan twee cijfers: uw \textbf{meest gewonnen
drankje} en het \textbf{totaal aantal consumpties dat u hebt getrakteerd}.

\rudi[rudi-gemeen]{Dat tweede getal is het eerlijkste klassement van de club.
Een rating praat u nog goed met ,,ongelukkige loting''. Zeventien getrakteerde
pinten niet. Dat is geen statistiek, dat is een bekentenis.}

\begin{center}
{\color{smaragd}\bfseries Proost --- en drink met mate. De winnaar heeft er
niets aan als u volgende week niet komt opdagen.}
\end{center}

\end{document}
